\documentclass[11pt,a4paper]{article}

\usepackage[margin=1in]{geometry}
\usepackage{amsmath,amssymb,booktabs,array,longtable,enumitem}
\usepackage{xcolor}
\usepackage{hyperref}

\hypersetup{
  colorlinks=true,
  linkcolor=blue!55!black,
  urlcolor=blue!55!black,
  citecolor=blue!55!black
}

\title{Research Plan: Zhang-Style HAR, GHAR, and GNNHAR-IV Across Dow 30, S\&P 100, and S\&P 500}
\author{Volatility Project}
\date{\today}

\begin{document}
\maketitle

\section{Objective}

The project will build a reproducible empirical study following the forecasting framework of Zhang Chao's GNNHAR work, then extend it with implied volatility. The study will compare Dow 30, S\&P 100, and S\&P 500 universes using each stock's close price, return, historical volatility, and implied volatility.

The central research questions are:
\[
\text{Does graph information improve volatility forecasts beyond HAR?}
\]
\[
\text{Does a nonlinear GNNHAR graph model improve beyond linear GHAR?}
\]
\[
\text{Does implied volatility add genuine predictive information beyond extra parameters?}
\]
\[
\text{Do the graph and IV gains change as the universe size grows from } N=30 \text{ to } N=100 \text{ and } N=500?
\]

The main scale-effect quantities will be
\[
\Delta^{GHAR}_{N}=L(HAR_N)-L(GHAR_N),
\]
\[
\Delta^{GNN}_{N,k}=L(GHAR_N)-L(GNNHAR{k}L_N),
\]
\[
\Delta^{IV}_{N,m}=L(m_N)-L(m\text{-}IV_N),
\]
where \(L(\cdot)\) is out-of-sample MSE or QLIKE and \(m\) is HAR, GHAR, or GNNHAR.

\section{Existing Data and Gaps}

The current repository already has most of the required data.

\begin{longtable}{llll}
\toprule
Universe & Main files & Approximate shape & Date range \\
\midrule
\endhead
Dow 30 & \texttt{experiments/dow30/data} & 30 assets, 1250 trading days & 2021-01-26 to 2026-01-26 \\
S\&P 100 & \texttt{data/scale\_experiment/sp100} & 101 assets, 1256 trading days & 2021-06-09 to 2026-06-09 \\
S\&P 500 & \texttt{data/scale\_experiment/sp500} & 503 assets, 1256 trading days & 2021-06-09 to 2026-06-09 \\
\bottomrule
\end{longtable}

Available panels:
\begin{itemize}[leftmargin=*]
  \item historical-volatility or realized-volatility target panel \(v_{i,t}\);
  \item implied-volatility panel \(x_{i,t}\);
  \item daily return panel \(r_{i,t}\), currently used to estimate the GLASSO graph.
\end{itemize}

One gap remains: close-price panels are not stored as first-class research inputs in the current repository. Returns exist, but the plan should fetch and save adjusted close panels so that the return construction is auditable:
\[
r_{i,t}=\log(P^{adj}_{i,t})-\log(P^{adj}_{i,t-1}).
\]
Close prices should not be used directly as model predictors in the main specification because price levels are nonstationary and not part of the Zhang HAR/GHAR/GNNHAR volatility setup. They are required for audit, return reconstruction, corporate-action checks, and robustness documentation.

\section{Data Design}

For each universe \(U_N\), build a synchronized panel with common dates and common tickers:
\[
\mathcal{D}_N=\{P^{adj}_{i,t}, r_{i,t}, v_{i,t}, x_{i,t}: i\in U_N, t=1,\ldots,T_N\}.
\]

The preferred target is the existing historical-volatility close-to-close 30-day panel when using the VolVue S\&P data. In the manuscript, this should be described carefully as the volatility target \(v_{i,t}\). It should not be overstated as high-frequency realized variance unless a true intraday realized-variance panel is later added.

The minimum data cleaning rules are:
\begin{itemize}[leftmargin=*]
  \item align close, return, historical-volatility, and IV panels by date and ticker;
  \item require at least 98 percent coverage in the main specification;
  \item fill short internal gaps only, for example with time interpolation capped at five trading days;
  \item winsorize or flag extreme volatility observations before final reporting;
  \item keep raw downloaded close prices, cleaned close prices, returns, volatility panels, and manifests.
\end{itemize}

Recommended output layout:
\begin{verbatim}
data/research_panels/dow30/
data/research_panels/sp100/
data/research_panels/sp500/
\end{verbatim}
Each folder should contain \texttt{adjusted\_close.csv}, \texttt{daily\_returns.csv}, \texttt{merged\_rv\_data\_filled.csv}, \texttt{merged\_iv\_data\_filled.csv}, \texttt{tickers.txt}, and \texttt{panel\_manifest.json}.

\section{Zhang-Style Feature Construction}

The base HAR feature vector for stock \(i\) at forecast date \(t\) is
\[
H^{RV}_{i,t-1}
=
\left(
v_{i,t-1},
\frac{1}{5}\sum_{\ell=1}^{5}v_{i,t-\ell},
\frac{1}{22}\sum_{\ell=1}^{22}v_{i,t-\ell}
\right).
\]

The IV extension uses the same daily, weekly, and monthly lag structure:
\[
H^{IV}_{i,t-1}
=
\left(
x_{i,t-1},
\frac{1}{5}\sum_{\ell=1}^{5}x_{i,t-\ell},
\frac{1}{22}\sum_{\ell=1}^{22}x_{i,t-\ell}
\right).
\]

The forecasting target is \(v_{i,t}\). The main model input for IV models is either
\[
X_{i,t-1}=H^{RV}_{i,t-1}
\]
or
\[
X^{IV}_{i,t-1}=\left[H^{RV}_{i,t-1},H^{IV}_{i,t-1}\right].
\]

\section{Graph Construction}

Following Zhang's setup, graph structure should come from the return panel, not from the volatility target. For each rolling forecast origin, estimate a graph from the training-window return matrix:
\[
R_t=(r_{1,t},\ldots,r_{N,t}).
\]

The main graph is the nonzero pattern of a GLASSO precision matrix:
\[
\widehat{\Theta}
=
\arg\max_{\Theta \succ 0}
\left\{
\log\det(\Theta)-\operatorname{tr}(\widehat{\Sigma}\Theta)-\lambda\|\Theta\|_1
\right\}.
\]
The adjacency matrix \(W\) is constructed from off-diagonal nonzero entries of \(\widehat{\Theta}\), symmetrized, and normalized.

Required graph robustness checks:
\begin{itemize}[leftmargin=*]
  \item GLASSO with cross-validated \(\lambda\);
  \item fixed-\(\lambda\) GLASSO for reproducibility;
  \item correlation graph fallback for numerical failures;
  \item density-matched random graph placebo;
  \item optional degree cap for S\&P 500 if the graph is too dense.
\end{itemize}

\section{Model Grid}

Every universe should be run with the same model family.

\begin{longtable}{lll}
\toprule
Group & Models & Main purpose \\
\midrule
\endhead
No graph, no IV & HAR & baseline volatility autoregression \\
Linear graph & GHAR & graph spillover gain beyond HAR \\
Nonlinear graph & GNNHAR1L, GNNHAR2L, GNNHAR3L, GNNHAR4L, GNNHAR5L & nonlinear and multi-hop graph gain \\
IV baseline & HAR-IV & option-market information gain without graph \\
Linear graph with IV & GHAR-IV & graph plus IV interaction in a linear model \\
GNN graph with IV & GNNHAR1L-IV through GNNHAR5L-IV & main extension \\
Fake-IV placebo & HAR-fakeIV, GHAR-fakeIV, GNNHARkL-fakeIV & parameter-expansion control \\
Graph placebo & GHAR-random, GNNHARkL-random & graph-structure control \\
\bottomrule
\end{longtable}

The current rolling pipeline already parses layer counts from model names like \texttt{GNNHAR4L}; therefore, 4L and 5L should be implemented through the newer Zhang-style pipeline rather than the older \texttt{src/gnnhar/models.py}, which only hard-codes 1L, 2L, and 3L.

\section{Linear Model Specifications}

The HAR model is
\[
v_{i,t}=\alpha_i+\beta_i^\top H^{RV}_{i,t-1}+\varepsilon_{i,t}.
\]

The GHAR model adds graph-aggregated volatility features:
\[
v_{i,t}
=
\alpha_i+\beta_i^\top H^{RV}_{i,t-1}
+\gamma_i^\top (W H^{RV}_{t-1})_i
+\varepsilon_{i,t}.
\]

The HAR-IV model is
\[
v_{i,t}
=
\alpha_i+\beta_i^\top H^{RV}_{i,t-1}
+\delta_i^\top H^{IV}_{i,t-1}
+\varepsilon_{i,t}.
\]

The GHAR-IV model is
\[
v_{i,t}
=
\alpha_i+\beta_i^\top H^{RV}_{i,t-1}
+\gamma_i^\top (W H^{RV}_{t-1})_i
+\delta_i^\top H^{IV}_{i,t-1}
+\eta_i^\top (W H^{IV}_{t-1})_i
+\varepsilon_{i,t}.
\]

As a robustness check, GHAR can also include higher graph powers \(W^2\) and \(W^3\), but the main Zhang-aligned comparison should keep GHAR as the linear graph baseline and interpret GNNHAR layers as the nonlinear multi-hop extension.

\section{GNNHAR Specifications}

The GNNHAR input at each date is a node-feature tensor. Without IV:
\[
X_{t-1}\in\mathbb{R}^{N\times 3}.
\]
With IV:
\[
X^{IV}_{t-1}\in\mathbb{R}^{N\times 6}.
\]

A \(k\)-layer GNNHAR model applies graph propagation and nonlinear transformations:
\[
Z^{(0)}=X_{t-1},
\]
\[
Z^{(\ell)}=\sigma\left(A Z^{(\ell-1)} B_{\ell}+Z^{(\ell-1)} C_{\ell}\right),
\quad \ell=1,\ldots,k,
\]
followed by a node-level forecast head producing \(\widehat{v}_{i,t}\).

The research grid will run \(k=1,2,3,4,5\). The 4L and 5L models are important for the scale hypothesis, but they should be interpreted carefully because deeper graph models can oversmooth and become unstable, especially for S\&P 500 and QLIKE training.

\section{Fake-IV and Placebo Design}

Fake IV should preserve the scale and cross-sectional shape of IV while breaking predictive timing. The default placebo is date permutation:
\[
\widetilde{H}^{IV}_{i,t}=H^{IV}_{i,\pi(t)}.
\]

The IV contribution can then be decomposed as:
\[
\text{total IV gain}=L(m)-L(m\text{-}IV),
\]
\[
\text{parameter gain}=L(m)-L(m\text{-}fakeIV),
\]
\[
\text{genuine IV information gain}=L(m\text{-}fakeIV)-L(m\text{-}IV).
\]

This decomposition is required for HAR, GHAR, and the best-performing GNNHAR depth.

\section{Rolling Forecast Protocol}

The main experiment should follow Zhang's rolling forecast structure:
\begin{itemize}[leftmargin=*]
  \item use up to 1000 prior observations for estimation;
  \item reserve the latest 22 pre-origin observations for validation in neural models;
  \item forecast the next 22 trading days as one block;
  \item move the forecast origin forward by 22 trading days;
  \item recompute GLASSO using only pre-origin returns for every block;
  \item never use future returns, future IV, or future volatility in graph construction or features.
\end{itemize}

For local smoke tests, run only one or two blocks. For final Colab runs, use all feasible blocks after the first origin. If runtime is too high for S\&P 500, use a staged final run: first MSE-trained models for all depths, then QLIKE-trained models for selected depths.

\section{Evaluation}

Each model should be evaluated out of sample with:
\[
MSE=\frac{1}{|\mathcal{T}|N}\sum_{t,i}(v_{i,t}-\widehat{v}_{i,t})^2,
\]
\[
QLIKE=\frac{1}{|\mathcal{T}|N}\sum_{t,i}
\left(
\log(\widehat{v}_{i,t})
+\frac{v_{i,t}}{\widehat{v}_{i,t}}
\right).
\]

The report should distinguish estimation loss from evaluation loss. A model trained under MSE should still be evaluated with both MSE and QLIKE.

Required tables:
\begin{itemize}[leftmargin=*]
  \item model loss table by universe, model, training criterion, MSE, and QLIKE;
  \item loss ratio table relative to HAR and relative to HAR-IV;
  \item MCS table at the 5 percent level;
  \item Diebold-Mariano pairwise comparison table;
  \item real-IV versus fake-IV decomposition table;
  \item regime-split table for calm and volatile periods;
  \item graph audit table with edge count, density, GLASSO alpha, and fallback flags.
\end{itemize}

Required figures:
\begin{itemize}[leftmargin=*]
  \item GLASSO adjacency heatmaps for each universe;
  \item loss-ratio bar charts across universes;
  \item forecast error boxplots;
  \item IV decomposition chart;
  \item GNN depth curve for 1L through 5L;
  \item graph density and fallback diagnostics.
\end{itemize}

\section{Statistical Tests}

Run MCS at the 5 percent level separately for each universe and evaluation loss. Run DM tests for the following pairs:
\begin{itemize}[leftmargin=*]
  \item HAR versus GHAR;
  \item GHAR versus each GNNHAR depth;
  \item HAR-IV versus GHAR-IV;
  \item GHAR-IV versus each GNNHAR-IV depth;
  \item each real-IV model versus its fake-IV counterpart;
  \item GNNHAR1L-IV versus GNNHAR2L-IV, and adjacent depth comparisons through 5L.
\end{itemize}

Because forecasts are block-based and cross-sectional, standard errors should be reported with care. At minimum, keep ticker-date loss records so that block bootstrap or date-clustered variants can be added.

\section{Colab Execution Plan}

No full Colab run should begin until the local data audit and smoke tests pass.

\subsection*{Stage 1: Local audit}
\begin{enumerate}[leftmargin=*]
  \item Create or refresh adjusted-close panels for Dow 30, S\&P 100, and S\&P 500.
  \item Recompute returns from close prices and compare to existing return panels.
  \item Build \texttt{panel\_manifest.json} for each universe.
  \item Verify ticker counts, date ranges, missingness, and units.
\end{enumerate}

\subsection*{Stage 2: Local smoke tests}
\begin{enumerate}[leftmargin=*]
  \item Run Dow 30 with HAR, GHAR, HAR-IV, GHAR-IV, GNNHAR1L, and GNNHAR1L-IV for one block.
  \item Run S\&P 100 with 1L through 5L for one block.
  \item Run S\&P 500 with HAR/GHAR and one GNN depth for one block.
  \item Confirm outputs are finite, positive, and written to the expected output folders.
\end{enumerate}

\subsection*{Stage 3: Colab final runs}
\begin{enumerate}[leftmargin=*]
  \item Upload or mount the audited panels in Google Drive.
  \item Clone the repository in Colab and install requirements.
  \item Run Dow 30 full rolling experiment.
  \item Run S\&P 100 full rolling experiment.
  \item Run S\&P 500 full rolling experiment, starting with MSE training and selected QLIKE training if time permits.
  \item Save all predictions, losses, graph audits, tables, figures, and logs to Drive.
\end{enumerate}

\subsection*{Stage 4: Report}
\begin{enumerate}[leftmargin=*]
  \item Summarize whether GHAR improves on HAR as \(N\) grows.
  \item Summarize whether GNNHAR improves on GHAR and which depth is best.
  \item Summarize whether IV gains are genuine using fake-IV controls.
  \item Discuss limitations: volatility-target definition, IV source quality, current-constituent bias, graph stability, and deep-GNN oversmoothing.
\end{enumerate}

\section{Proposed Commands}

The final commands should use the existing Zhang-style rolling pipeline. The exact paths can be adjusted after the audited panels are written.

\begin{verbatim}
python scripts/analysis/gnnhar_iv_zhang_scale_pipeline.py \
  --universe-name sp100 \
  --data-dir data/research_panels/sp100 \
  --returns-file data/research_panels/sp100/daily_returns.csv \
  --output-dir outputs/final-zhang-scale/sp100 \
  --models HAR,GHAR,HAR+IV,GHAR+IV,GNNHAR1L,GNNHAR2L,GNNHAR3L,GNNHAR4L,GNNHAR5L,GNNHAR1L-IV,GNNHAR2L-IV,GNNHAR3L-IV,GNNHAR4L-IV,GNNHAR5L-IV \
  --lookback 1000 \
  --window 22 \
  --valid-len 22 \
  --graph-method glasso_cv \
  --loss MSE
\end{verbatim}

Run analogous commands for \texttt{dow30} and \texttt{sp500}. Add fake-IV and random-graph placebo models after the first full baseline run is verified.

\section{Acceptance Criteria}

The plan is complete only when the final artifact contains:
\begin{itemize}[leftmargin=*]
  \item audited data manifests for Dow 30, S\&P 100, and S\&P 500;
  \item full rolling predictions for all requested model families;
  \item MSE and QLIKE evaluation for every model;
  \item MCS and DM statistical comparisons;
  \item real-IV versus fake-IV decomposition;
  \item graph diagnostics for every rolling block;
  \item a final PDF or LaTeX report that clearly separates Zhang's original framework from the IV extension.
\end{itemize}

\end{document}
